% =====================================================================
%  二、问题分析
%  写法约定：每问只说明三件事——考虑到了什么、要建立什么样的模型、
%  为什么最终采用这一做法；推导细节留给第五章。
% =====================================================================
\section{问题分析}

四个子问题层层递进。问题一是单次交会定位的几何基础；问题二则将该几何模型进一步用作
检测点位置的选择依据，形成单源定位策略；问题三在严格的时间预算下，把单源定位能力嵌入“测向、收敛、
清除”的在线决策；问题四基于问题三模型的基础之上，再额外引入定向源方向诊断对决策的影响。

\subsection{问题一的分析}
本问不宜假设误差分布，而应利用题目给出的误差界，采用最坏情形模型：
每个检测点的示向度对应一个由“测量方向 \(\pm 1^\circ\)”张成的角楔。
多个检测点的角楔求交，即得干扰源的可能定位区域。由于角楔是凸集，其交仍为凸多边形，
因此区域直径可转化为凸多边形顶点间的最远点对问题。求解时，先判断交会区域是否有界，
再检验以该直径为直径的圆能否覆盖定位区域。整体求解思路可概括为：由误差界构造单点角楔，
通过多检测点角楔交会确定凸多边形定位区域，进而将区域直径计算归结为顶点间的最远点对问题；
同时辅以有界性判定与圆覆盖检验。

\subsection{问题二的分析}

\textrm{【待填充：确定该准则并据此给出选择策略与候选区域。】}

\subsection{问题三的分析}

\textrm{【待填充：给出策略与算法，并补充模拟与实测指标。】}

\subsection{问题四的分析}

\textrm{【待填充：给出定向源的检测与清除策略、诊断规则及实测结果。】}
